\begin{frame}{ELBO: 계산 가능한 하한으로 우회}
  직접 계산할 수 없는 $\log p_\theta(x)$ 대신, 다음이 성립함을 보일 수
  있다:
  \[
  \log p_\theta(x) \ge
  \mathbb{E}_{z \sim q_\phi(z|x)}[\log p_\theta(x|z)]
  - D_{KL}\big(q_\phi(z|x) \,\|\, p(z)\big)
  \]
  우변을 \kb{ELBO}(Evidence Lower BOund)라 부른다. 두 항의 의미:
  \vskip0.6em
  \begin{block}{첫 항: 복원 항}
    $\mathbb{E}[\log p_\theta(x|z)]$ — 인코더가 만든 $z$로 디코더가
    원본 $x$를 얼마나 잘 복원하는가 (일반 오토인코더의 복원 오차와
    본질적으로 같음).
  \end{block}
  \vskip0.4em
  \begin{block}{둘째 항: 정규화 항}
    $D_{KL}(q_\phi(z|x) \| p(z))$ — 인코더 분포가 목표 사전분포 $p(z)$
    (보통 표준정규)에서 얼마나 벗어났는가 (KL divergence, Week 2.3
    섀넌의 정보이론에서 소개). 이 항이 바로 잠재 공간을 매끄럽게
    만드는 힘.
  \end{block}
\end{frame}
